
\subsection[Category of Instantiations]{The category of instantiations $\instC$}

The category $\instC$ is defined as
\begin{itemize}
    \item Objects: worlds $\vw = \{ \vl_1 : \vC_1, \dots, \vl_\vn : \vC_\vn \}$.
    \item Morphisms: Injective, sort-preserving $\rho : \vw \to \vw'$ and instantiation data $\mvarM : \dheapF(\vw' \ominus \rho, \vw')$.
\end{itemize}
We have a forgetful functor $\instU : \instC \to \worldC$ that forgets the instantiation data.

\begin{definition}[Inhabited Sorts]\label{def:inhabited-sorts}
    A sort $\vC \in \typeSort$ is said to be \emph{inhabited} if
    there exists a world $\vw$, renaming $\va : \vC \to \vw$, and $\mvarM : \heapF(\vw)$.
\end{definition}
Similarly we say that a set of sorts is inhabited if every sort in the set is inhabited.

\begin{lemma}\label{lem:lic-instC}
    The local independent structure on $\worldC$ can be directly extended to $\instC$ such that $\instU$ reflects local independent squares.
\end{lemma}

Recall that the hiding monad $\hideMon$ on $\instC$ is given by the colimit $\hideMon\vA\vw = \Colimit{(\sliceC\vw\instU \xto\vU \instC \xto\vA \setC)}$.
The following lemma describes the quotienting that happens by this colimit.
\begin{lemma}[Quotienting in the Hiding Monad]\label{lem:hmon-quotient}
    For $\sigma_1 : \vw \to \vw_1$, $\va_1 \in \vA(\vw_1)$, and $\sigma_2 : \vw \to \vw_2$, $\va_2 \in \vA(\vw_2)$, 
    we have that $\eqclass{\sigma_1}\va_1 = \eqclass{\sigma_2}\va_2 \in \hideMon\vA\vw$ iff there exist 
    \begin{tikzcd}
        & {\hat\vw}
    \\ 
    {\vw_1} 
        \ar[ur, "\tinstU\me_1'"] 
            && {\vw_2} 
                \ar[ul, "\tinstU\me_2'"']
    \\ 
        & {\vw} 
            \ar[ul, "\sigma_1"] \ar[ur, "\sigma_2"']
    \end{tikzcd}
    s.t. $\vA(\me_1')(\va_1) = \vA(\me_2')(\va_2)$.
\end{lemma}
\begin{proof}
    We first show that the relation
    \[ \inj_{\sigma_1}(\va_1) \sim \inj_{\sigma_2}(\va_2) \iff 
       \ExistsTwo{\me_1'}{\me_2'} \instU\me_1' \comp \sigma_1 = \instU\me_2' \comp \sigma_2 \Land \vA(\me_1')(\va_1) = \vA(\me_2')(\va_2) \]
    is an equivalence.
    Reflexivity and symmetry is immediate.
    For transitivity, we have
    \begin{center}
        \begin{tikzcd}
                && {\vw_{12} \liplus{\vw_2} \vw_{23}}
        \\
            & {\vw_{12}} 
                \ar[ur, "\ext^{\me_{22}}\me_{23}"]
                    && {\vw_{23}}
                        \ar[lu, "\ext^{\me_{23}}\me_{22}"']
        \\
        {\vw_1}
            \ar[ur, "\tinstU\me_1"]
                && {\vw_2}
                        \ar[ul, "\tinstU\me_{22}"] \ar[ur, "\tinstU\me_{23}"']
                        && {\vw_3}
                            \ar[ul, "\tinstU\me_3"']
        \\\\
                && {\vw_1}
                    \ar[uull, "\sigma_1"] \ar[uu, "\sigma_2"] \ar[uurr, "\sigma_3"']
        \end{tikzcd}
    \end{center}
    Notice that the top square commutes by \cref{lem:lic-instC}.
    which implies that $\inj_{\sigma_1}(\va_1) \sim \inj_{\sigma_3}(\va_3)$.

    To prove that the generated relation by the colimit is exactly $\sim$, one shows direction $\to$ by inducting on the list of related elements,
    and for $\from$, notice that if $\inj_{\sigma_1}(\va_1) \sim \inj_{\sigma_3}(\va_3)$, then
    \[ \inj_{\sigma_1}\va_1 \leq \inj_{\instU\me_1 \comp \sigma_1}\vA(\me_1)(\va_1) \ge \inj_{\sigma_2}\va_2 \]

\end{proof}
